\documentclass[../../../../main.tex]{subfiles}
\begin{document}

\begin{flushright}
\footnotesize\textit{【自動文字起こし・要確認】}
\end{flushright}

\noindent
\textbf{[解]}

\begin{center}
\begin{tikzpicture}[scale=0.8, >=stealth]
\draw[->] (-0.5,0) -- (4.5,0) node[right] {$x$};
\draw (0.8, 0.8) circle (0.8);
\draw (2.5, 1.2) circle (1.2);
\draw (1.55, 0.25) circle (0.25);
\node at (0.8, 0.8) {$C_{n+1}$};
\node at (2.5, 1.2) {$C_n$};
\node at (1.55, 0.45) {$C_{n+2}$};
\filldraw (0.8,0) circle (1pt) node[below] {$x_{n+1}$};
\filldraw (2.5,0) circle (1pt) node[below] {$x_n$};
\filldraw (1.55,0) circle (1pt) node[below] {$x_{n+2}$};
\node[above] at (2, 2.5) {$n$ : odd の時};
\end{tikzpicture}
\qquad
\begin{tikzpicture}[scale=0.8, >=stealth]
\draw[->] (-0.5,0) -- (4.5,0) node[right] {$x$};
\draw (0.8, 1.2) circle (1.2);
\draw (2.5, 0.8) circle (0.8);
\draw (1.75, 0.25) circle (0.25);
\node at (0.8, 1.2) {$C_n$};
\node at (2.5, 0.8) {$C_{n+1}$};
\node at (1.75, 0.45) {$C_{n+2}$};
\filldraw (0.8,0) circle (1pt) node[below] {$x_n$};
\filldraw (2.5,0) circle (1pt) node[below] {$x_{n+1}$};
\filldraw (1.75,0) circle (1pt) node[below] {$x_{n+2}$};
\node[above] at (2, 2.5) {$n$ : even の時};
\end{tikzpicture}
\end{center}

円 $C_n$ の中心 $O_n$ とする. $O_{n+2}$ から $x = x_n, x_{n+1}$ への垂足を各 $H_n, H_{n+1}$ とおく.
$\triangle O_n H_n O_{n+2}, \triangle O_{n+1} H_{n+1} O_{n+2}$ に三平方を用いて,
\[
\begin{cases}
(r_{n+2} + r_n)^2 = (x_n - x_{n+2})^2 + (r_{n+2} - r_n)^2 \\
(r_{n+2} + r_{n+1})^2 = (x_{n+1} - x_{n+2})^2 + (r_{n+2} - r_{n+1})^2
\end{cases}
\]
\[
\begin{cases}
2\sqrt{r_n r_{n+2}} = |x_n - x_{n+2}| \\
2\sqrt{r_{n+1} r_{n+2}} = |x_{n+2} - x_{n+1}|
\end{cases}
\quad (\because r_k > 0)
\]
だから, $n$: even の時 $x_n < x_{n+2} < x_{n+1}$, $n$: odd の時 $x_{n+1} < x_{n+2} < x_n$ であることより,
\[
\begin{cases}
2\sqrt{r_n r_{n+2}} = (-1)^n (x_{n+2} - x_n) \\
2\sqrt{r_{n+1} r_{n+2}} = (-1)^n (x_{n+1} - x_{n+2})
\end{cases}
\quad \dots \text{①}
\]
又, $O_n, O_{n+1}$ に対しても同様に
\[
(r_{n+1} + r_n)^2 = (r_{n+1} - r_n)^2 + (x_{n+1} - x_n)^2
\]
\[
2\sqrt{r_n r_{n+1}} = (-1)^n (x_{n+1} - x_n) \quad \dots \text{②}
\]
が成立する.

\bigskip

(1) $q_n \in \mathbb{Z} \dots \spadesuit$ を帰納法で示す. $n=0, 1$ の時, $q_0 = q_1 = 1$ から $\spadesuit$ は成立するので, 以下 $n=k, k+1 \ (k \in \mathbb{Z}_{\ge 0})$ での成立を仮定する.
①で辺々足したものと②から
\[
2\sqrt{r_k r_{k+2}} + 2\sqrt{r_{k+1} r_{k+2}} = 2\sqrt{r_k r_{k+1}}
\]
$q_n = \frac{1}{2\sqrt{r_n}}$ から
\[
\frac{1}{q_k q_{k+2}} + \frac{1}{q_{k+1} q_{k+2}} = \frac{1}{q_k q_{k+1}}
\]
両辺に $q_k q_{k+1} q_{k+2}$ をかけて,
\[
q_{k+2} = q_k + q_{k+1} \in \mathbb{Z} \quad (\because \text{仮定}) \quad \dots \text{③}
\]
よって $n=k+2$ でも $\spadesuit$ は成立. 以上から $\spadesuit$ は示された. \hfill $\blacksquare$

\bigskip

(2) ①の辺々わって, $p_n = \frac{x_n}{2\sqrt{r_n}}$ から
\[
\begin{cases}
(-1)^n = (q_n p_{n+2} - q_{n+2} p_n) & \dots \text{④} \\
(-1)^n = (q_{n+2} p_{n+1} - q_{n+1} p_{n+2}) & \dots \text{⑤}
\end{cases}
\]
④, ⑤を代入して引いて,
\[
q_n(p_{n+1} + p_n) + q_{n+1}(p_n + p_{n+1}) = p_{n+2}(q_n + q_{n+1})
\]
$q_n + q_{n+1} \neq 0$ だから ($\because q_n > 0$),
\[
p_n + p_{n+1} = p_{n+2} \quad \dots \text{⑥}
\]
が成り立つ. 以下, $p_n \in \mathbb{Z}$ かつ $p_n$ と $q_n$ が互いに素 $\dots \clubsuit$ であることを帰納的に示す.
$n=0, 1$ の時, $p_0 = 0, p_1 = 1$ となって $\clubsuit$ は成立するので, 以下 $n=k, k+1 \ (k \in \mathbb{Z}_{\ge 0})$ での成立を仮定する.
すると⑥から, $p_{k+2} \in \mathbb{Z}$ となる $\dots \text{⑦}$. さらに, 初期条件及び③, ⑥から,
$p_{n+1} = q_n$ となるので, 「$p_{k+2}$ と $q_{k+2}$ が互いに素」$\iff$「$q_{k+1}$ と $p_{k+2}$ が互いに素」$\iff$「$q_{k+1}$ と $q_{k+1} + q_k$ が互いに素」$(\because \text{③}) \iff$「$q_{k+1}$ と $q_k$ が互いに素」$\iff$「$q_{k+1}$ と $p_{k+1}$ が互いに素」となり, 仮定から $q_{k+2} \perp p_{k+2}$ となる $\dots \text{⑧}$
⑦, ⑧から $n=k+2$ でも $\clubsuit$ は成立. よって $\clubsuit$ は示された. \hfill $\blacksquare$

\bigskip

(3) (2)から, $p_{n+1} = q_n \ (n \ge 0)$. $p_0 = 0$ であるので,
\[
x_n = \begin{cases} 0 & (n = 0) \\ \dfrac{q_{n-1}}{q_n} & (n \ge 1) \end{cases}
\]
となる. ③の両辺 $q_k$ でわって,
\[
\frac{1}{x_{n+2}} = x_{n+1} + 1 \quad (n \ge 0)
\]
これより, $x_{n+1} = \dfrac{1}{1 + x_n} \quad (n \ge 1)$.
両辺 $\alpha = \dfrac{1}{1 + \alpha}$ を引いて,
\[
x_{n+1} - \alpha = \frac{1}{1 + x_n} - \frac{1}{1 + \alpha} = \frac{-1}{(1 + \alpha)(1 + x_n)} (x_n - \alpha)
\]
両辺絶対値をとって,
\[
|x_{n+1} - \alpha| = \left| \frac{1}{(1 + \alpha)(1 + x_n)} \right| |x_n - \alpha|
\]
ここで, $\alpha = \dfrac{-1 + \sqrt{5}}{2}$, $0 \le x_n \le 1$ から
\[
|x_{n+1} - \alpha| \le \left| \frac{\sqrt{5}-1}{2} \right| |x_n - \alpha| < \frac{2}{3} |x_n - \alpha|
\]
なので, くり返し用いて,
\[
|x_n - \alpha| < \left(\frac{2}{3}\right)^{n-2} |x_2 - \alpha|
\]
はさみうちから, $x_n \longrightarrow \alpha = \dfrac{-1 + \sqrt{5}}{2}$ である.

\newpage

\end{document}
